\section{齐：富庶东方与稷下的另一条国家道路}

从临淄向东，陆路穿过半岛，水道、盐泽和可停泊转运的海滨把齐接到大海；向西，则是黄河下游与中原诸侯往来的道路。沿岸港湾并不自动等于一支军队，盐铁也不自动等于一套无所不能的财政机器，可它们同城邑、田土和都市市场相接，能不断汇来粮食、器物、货币与人。战国的齐因此不只是一支军队或一座都城：它有可供征发的城邑，有愿意为俸禄和名声而来的游士，也有一旦失去盟友便格外显眼的财富。

田氏取得齐国以后，正是在这块富庶的底盘上，把新的家族统治、列国外交和文化声望接到一处。它后来的崩解也从这里显出代价：能汇聚资源的临淄，在战争中也会成为最醒目的目标；能够越过海滨和中原道路流动的货物与人，并不能替国家守住盟约、渡口和指挥链。

\subsection{易姓之后：把东方的资源变成国家能力}

田氏不是在前386年才突然拥有齐国。旧姜姓公室衰微的过程中，田氏已同邑宰、宗族和临淄周边的地方关系相互嵌合。关于这种基层结盟，传世文献只留下由上而下的家族政治叙事，墓葬和聚落材料也难让我们点出某一村社在何年归附；不过，若没有地方精英、劳力和粮食的持续合作，册命便只是都城里的一纸仪式。这个较长的积累过程，可由\eventanchor{WQ-0428-TIANQI-RURAL-ALLIANCE}{战国·田氏地方结盟}来理解，而不能压缩成一夜夺门的故事。

同样应当谨慎看待齐的“富”。\eventanchor{WQ-0394-QI-TIAN-REVENUE}{战国·临淄财赋}所指的，是临淄周边盐、铁、城邑赋入与市场网络逐步被动员的趋势，并非保存下来的田齐财政决算。《管子》有关轻重、盐铁和物价的篇章成书层次复杂，不能移作前四世纪某年颁下的财政令；临淄故城及铸造遗址则能提示都市生产和物资集散的重要性，不能替我们算出国库的每一笔收入。较稳妥的说法是，东向海滨、半岛道路和临淄市场给田氏提供了更宽的征发与赏赐余地，也使调粮、燃料、劳役和守备的安排变得更加繁复。

\begin{event}{公元前386年}{田氏获周册命，齐国易姓}{可靠纪年}{临淄与周王室}\eventanchor{WQ-0386-TIANQI}{战国·田氏代齐}
田氏取代姜氏，并非一日之间城门易帜。《史记·田敬仲完世家》把田和安置齐康公、向周廷求封与获命写在同一条权力转换的脉络中；《竹书纪年》的相关记载又提示旧君被迁置、隔离。迁往海上的具体地点、待遇和先后并不一致，故不能把它写成可逐项核对的押送路线；可以确认的是，\eventanchor{WQ-0386-TIANHE-EXILE-KANG}{战国·迁置齐康公}旨在把旧君从可能重新凝聚宗族与礼仪支持的中心移开。

到前386年，周王室承认田氏为齐侯，旧姜姓国君已不能支配齐国的军政资源。随后有关宗庙、爵赏和使聘的调整，见于后世叙事和田齐兵器铭文可以相互照亮的范围；\eventanchor{WQ-0386-TIANQI-RITUALS}{战国·田齐重定名分}并不意味着我们拥有一份完整的礼制清单。临淄仍是那座临淄，城邑、田土与通往海滨的利路也没有随姓氏改变；改变的是谁能以齐侯之名调度它们。

这层名分仍值得争取。它和\eventref{WQ-0403-THREEJIN}{战国·三家分晋}所见的情形相似：周廷保有授名的仪式权威，却无力逆转列国中已经形成的强弱。周命为田氏已取得的优势添上可以通行于列国之间的名号，也让新王室能把赏赐、使聘和征发说成“齐”的事务，而不只是强族私门的胜利。前379年姜齐康公去世，旧王族最后的政治象征进一步消退；\eventanchor{WQ-0379-QI-KANGGONG}{战国·姜齐终结}与\eventanchor{WQ-0379-QI-KANG-SHRINE}{战国·旧君宗祀}之间宗祀、封邑如何处置，材料并不充分，却显示易姓之后仍须处理礼制余波。
\end{event}

田齐要使新名分有效，不能只倚临淄的宫门。后世把齐威王与邹忌的故事写成一面能照见朝廷的镜子：君主因近臣、妻妾和左右的称赞而警觉，进而广开进谏之路。它见于《战国策·齐策》并经《史记》编排，是劝谏文本，而不是逐字保存的朝议记录；\eventanchor{WQ-0371-QI-WEIWANG-COURT}{战国·齐威王纳谏}所能把握的，是田齐以整饬吏治、疏通信息来塑造政治形象的记忆，不能据此复原一套固定的考课程序。

即墨大夫与阿大夫受赏罚的叙事也有同样的性质。\eventanchor{WQ-0357-QI-WEIWANG-RECRUITMENT}{战国·即墨阿地考课}把远郡守令、地方毁誉和君主判断压缩成鲜明对照，说明王权关心地方官是否截留民情、夸饰政绩，却不能证明齐廷每年都按这一故事的办法考核。西部山地的齐长城亦应这样读：考古调查可确认若干墙体、关隘与道路的长期防御意义，始建、续筑和延展到何处仍有争议。它作为\eventanchor{WQ-0360-QI-WALL}{战国·齐长城}，减少了陆路突袭的机会，也把维修、驻守和粮运的成本留给沿线城邑。

这些安排让齐能把主力投向外线。前354年赵向齐求援，齐廷不必立即在赵地同魏军硬拼，而要同时权衡临淄防务、援赵的代价和魏军能否被牵动。传世材料将田忌、孙膑安排在这一机动作战的中心；孙膑的生平、职权和兵法文本归属均有争议，仍可把\eventanchor{WQ-0354-QI-MILITARY-COMMAND}{战国·齐援赵军}视为齐国具有选择战场能力的证据。前342年马陵之役的“减灶”尤其著名，\eventanchor{WQ-0342-QI-REDUCTION-FIRES}{战国·马陵减灶}所见的孙膑形象，来自《史记·孙子吴起列传》与后起兵书、列传共同塑成的战术叙事；它或许保存了诱使魏军脱离稳固行军节奏的思路，却不能被当作现场军令的逐字记录。

齐在中原的地位由此不只是“东方富国”。它可以在赵、魏相争时选择救援，利用对手兵力与道路的分离来争取主动；这是一种外交位置，而非对其他国家的永久支配。齐军的机动、临淄的物资和田氏宫廷的声望互为条件，但每一次干预都会使被牵动的国家重新估算威胁。后来北向入燕的抉择，正是在这一能力与这一误判之间发生。

\subsection{稷下与干预：声望能打开道路，不能替代秩序}

前四世纪中叶以后，齐廷向学者提供居处、名号和议论空间的惯例逐步可见。把这概括为\eventanchor{WQ-0347-QI-JIXIA-PATRONAGE}{战国·稷下延士}与\eventanchor{WQ-0320-JIXIA}{战国·临淄稷下}，并不是说某一年忽然建立了一所具有校舍、学籍和统一课程的“学宫”。《史记·孟子荀卿列传》《史记·田敬仲完世家》与《盐铁论·论儒》保存的是不同年代、不同文体的回望；它们共同支持临淄曾以列第、俸禄和政治声望聚集游士，却不足以列出常设人数、经费账目或每日制度。

齐威王去世、齐宣王在前320年继位时，继承的正是这种能把财富转成名望的条件。\eventanchor{WQ-0320-QI-XUAN}{战国·齐宣王继位}之后，\eventanchor{WQ-0320-QI-XUAN-JIXIA}{战国·宣王延士}的材料所显示的不是单向授课，而是君主、权贵、门客和游士彼此利用的网络：齐廷希望获得治国、外交、军事和礼义的语言，游士则在临淄找到比单一宗族更宽的出路。孟子游齐的对话能使我们听见王道与功利相互争论的声音，却不是宫廷会议纪要；荀子与齐的关联、邹衍和淳于髡等姓名在后世文献中的聚集，也不能推成所有人同时受聘于同一机构。

\begin{textwindow}[《史记》笔下的稷下]
《史记·田敬仲完世家》说稷下诸士“皆赐列第为上大夫，不治而议论”。汉代史家把它置入齐国盛衰的叙事，说明后世记忆中的齐廷曾以居处、俸禄和议论聚士；它不能替我们把稷下画成制度完备、面貌不变的近代学院。
\end{textwindow}

稷下因而是一种国家能力，却是有限的国家能力。临淄的手工业与铸造作坊在前四世纪末至前三级初的扩展，见于故城与相关遗址；\eventanchor{WQ-0300-QI-LINZI-CRAFT}{战国·临淄手工业}表明都市对兵器、货币和日用品的集中需求可以支撑人才和外交的市场，也说明粮食、燃料、劳役和治安须由更复杂的调配来负担。知识网络让齐在诸侯之间显得富足、开放、可交往，却不能自动生产一个可靠的盟友集团。名士可以帮助国家说明何为善政，不能替它守住济西的渡口，更不能抹去一次外来干预在邻国留下的记忆。

\begin{event}{公元前314年}{齐军入燕，胜利留下北方的创伤}{可靠纪年}{燕地}\eventanchor{WQ-0314-QI-YAN}{战国·齐伐燕}
燕王哙、子之与太子一系的冲突撕裂了燕的命令链条，齐宣王遂遣兵北上。内乱中的道路、城邑和粮食供给本已失序，齐军得以深入；齐廷面对的选择却不只是“能不能打下”，而是打下之后由谁供粮、谁受命、谁愿意服从。外军的到来没有给燕地带来一套被当地精英和列国共同承认的新秩序。

燕人抵抗，列国亦不愿看见齐把北方危机化作常久领土。齐军终究撤离，\eventanchor{WQ-0314-QI-YAN-WITHDRAW}{战国·齐军撤燕}与\eventanchor{WQ-0314-QI-YAN-WITHDRAWAL-ORDER}{战国·齐军撤退}所指的过程并没有保存下命令文本；不能凭后世叙事给它补上一道精确的撤军诏令。更可把握的是，齐未能将军事进入转成稳定的占领和行政安排。随后围绕燕王位与占领政策的质疑，构成\eventanchor{WQ-0314-QI-YAN-LEGITIMACY}{战国·伐燕失信}的外交负担。

《孟子·公孙丑下》保存“齐能否取燕”的辩论，重点在王道和正当性；《史记·燕召公世家》则把它纳入燕由乱而复的国别叙事。两者都不能替我们精确还原每支部队的行程，正因为如此，它们反而清楚提出了干预的政治难题：以救乱或伐乱为名进入邻国，可以取得一时主动，却不能省去当地服从、列国承认和撤军之后的安排。前314年的损伤成为\eventref{WQ-YAN-ZHAO}{战国·燕昭王蓄力}得以凝聚人心的资源，也是后来\eventref{WQ-0284-QI}{战国·燕破齐}的近因之一。
\end{event}

齐在燕地留下的并非唯一压力。前301年宣王去世，湣王继位；\eventanchor{WQ-0301-QI-MIN-ASCEND}{战国·齐湣王继位}发生在齐对外影响仍强的时刻，也让下一代君主更可能把干预视为可重复使用的工具。齐、秦之间的人才流动尤其显示，外交并不只发生在国君之间。孟尝君田文入秦为相，\eventanchor{WQ-0298-MENGCHANG-QIN-CHANCELLOR}{战国·孟尝君相秦}的任相起讫已有异说；他在秦受疑而返齐的年份也难钉死。所谓“鸡鸣狗盗”把门客的机巧压缩成一段惊险的出关故事，见于《史记·孟尝君列传》和《战国策·齐策》，文学色彩极强，不能拿来复原函谷关的夜间程序。

但传记的夸张并不使问题失真。一个齐相在秦廷的去留，确会牵动两国对情报、人才和相权的猜疑；\eventanchor{WQ-0298-MENGCHANG-ESCAPE}{战国·孟尝君离秦}之后，田文返回齐国，依托门客和相国网络参与对秦外交，便是\eventanchor{WQ-0297-QI-MENGCHANG-RETURN}{战国·孟尝君返齐}所保留的结构性事实。至于他如何越过函谷关，仍应停留在\eventanchor{WQ-0297-MENGCHANG-HANGU-PASS}{战国·孟尝君出函谷}的传奇层次，而不是当作可核的边关档案。

前288年前后，齐湣王一度接受“东帝”称号，试图与秦的“西帝”相匹敌。称号的仪式、朝内反应和废除的确切节奏均不清楚；《史记》与《战国策》至少让我们看见，\eventanchor{WQ-0288-QI-EASTERN-EMPEROR-COURT}{战国·齐称东帝}是一种争取名分优势的试探，而迅速不再坚持该号，恰说明齐无法独自承担列国警惕。宫廷中的对秦路线也非铁板一块。吕礼在齐的相位和活动次序主要依《战国策》说辞重建，\eventanchor{WQ-0288-LULI-QI-CHANCELLOR}{战国·吕礼在齐}不宜被写成连续无缺的派系年表；相权竞争和强硬、审慎两种选择彼此牵制，则足以解释齐在扩张高峰时的外交并不稳定。

\begin{event}{公元前286年}{齐灭宋，扩张把财富变成共同的警报}{可靠纪年}{宋地、陶与彭城}\eventanchor{WQ-0286-QI-SONG}{战国·齐灭宋}
宋康王偃败亡，齐军取得宋地。齐所追求的，不只是地图上多一片颜色：泗水交通、陶与彭城一带的城邑、市场和人口，可能使齐更深入地进入中原。关于宋王的结局、齐军占领到何处，史书版本并不完全一致；《史记·宋微子世家》和《史记·田敬仲完世家》能够共同确认的，是前286年的政治归属发生了剧烈改变。

接收土地却比取得城门更难。我们没有齐廷接管陶、彭城赋役、货币和市场的成文档案，故\eventanchor{WQ-0286-QI-SONG-ADMINISTRATION}{战国·接收宋地}只能说齐尝试把这些资源纳入自身，而不能虚构郡县名称、税额或整齐的行政图。灭宋后的王偃形象同样带有后世道德化叙述，\eventanchor{WQ-0286-QI-SONG-WANGYAN}{战国·宋王偃败亡}不应由此延伸出未经证实的战场细节。

直接后果却十分实际：齐的新财富和新通道，也使赵、秦、韩、魏、燕更容易把它视为吞并威胁。齐曾以富强取得干预的余地，至此又因扩张超出外交承受力，成为各国可以共同防备的对象。燕昭王把雪耻、人材和对外联络慢慢编入国家准备，齐则一步步失去可用的回旋空间。
\end{event}

\subsection{济西之后：占领、复国与缩小的选择}

\begin{event}{公元前284年}{五国军渡济西，齐的防线崩解}{可靠纪年}{济西、临淄及齐地}\eventanchor{WQ-0284-QI}{战国·燕破齐}
乐毅统率燕、赵、秦、韩、魏五国之军攻齐。《史记·乐毅列传》所记济西之战，使齐军主力首先瓦解；盟军随后深入，临淄失守，齐湣王出奔，齐地大部落入燕军控制。传世叙事说乐毅下齐七十余城、改置燕郡，城数、设置范围和各地降附的节奏都不宜据此写得过于齐整。可确知的是，战败已从边境渡口延伸到临淄和齐国腹心。

临淄陷落意味着的也不只是宫室易主。它是田氏调度税入、工匠、使者和稷下声望的中心；\eventanchor{WQ-0284-YAN-CAPTURE-LINZI}{战国·燕军入临淄}之后，原来以都城为轴的行政和知识网络不得不依附尚存的城邑。关于杀掠和财物的数字，多来自胜利者或后出的国别叙事，应保留夸张的可能，不能以数字替代对占领后果的理解。

湣王自临淄东走，传世文献让他先后与卫、邹、鲁等地发生联系，路线与接待细节却彼此矛盾。应当保留\eventanchor{WQ-0284-QI-MINWANG-FLEES}{战国·齐湣王出奔}所显示的实质：王室一旦没有固定都城，旧盟友和地方精英便会重新计算忠属，而不是自动提供可持续的粮运和兵力。莒与即墨之所以重要，也不在于我们能列出城内难民名册；\eventanchor{WQ-0284-QI-JIMO-REFUGE}{战国·莒即墨守齐}所指的是士卒、宗室、粮食和地方组织仍可聚集于两座残城，为长期抵抗留下节点。

这不是燕国单独击倒一个富国的故事。燕为盟主而乐毅为将，是多年筹备的结果；赵、秦、韩、魏同时出兵，才是使能条件。各国未必共享同一种报复，却能共同利用齐孤立的时刻。齐的粮食、城邑和市场没有消失，反而成为攻取后的目标；它缺少的是让财富在危机中转化为外援、守备与可靠指挥的外交和军事关系。田氏齐由此失去战前能够左右列国局势的主动权。
\end{event}

流亡没有因王室还在就变得稳固。前283年，湣王在莒外为楚将淖齿所杀；《史记·田敬仲完世家》《史记·楚世家》叙述淖齿的动机和当地局势各有重心，\eventanchor{WQ-0283-QI-MINWANG-MURDER}{战国·齐湣王遇害}只能确定这次流亡王权再遭断裂，不能把任何一家的责难写成铁案。残余精英须以襄王和守城者重建合法性。关于莒地重建朝廷、收拢旧臣和粮道的地点、人事亦难逐项核实，\eventanchor{WQ-0283-QI-XIANG}{战国·齐襄王继位}与\eventanchor{WQ-0283-QI-XIANG-RESTORATION}{战国·莒地重建朝廷}标出的是这个必要而脆弱的政治中心。

人事转折出现在即墨。田单由低阶军吏升为统帅的过程主要见于列传，职位称谓不可过度制度化；不过，一座被围的城确实需要能够协调城防、工匠、乡里和守军的人。于是\eventanchor{WQ-0284-TIANDAN-JIMO}{战国·田单守即墨}与\eventanchor{WQ-0283-TIANDAN-JIMO-COMMAND}{战国·田单统即墨}不是一位英雄忽然出现的神话起点，而是残余政权把分散人力转为指挥能力的时刻。燕军方面也有裂缝：昭王去世后，燕惠王以骑劫代乐毅，乐毅奔赵。反间如何传递、惠王如何受说的铺陈多属《史记·田单列传》等后出叙事的戏剧化部分；易将、主将离去和占领军指挥受损的关联，则不能略去。

\begin{event}{约公元前279年至前278年}{田单组织反攻，齐地渐次收复}{年代细节有争议}{即墨及齐地}\eventanchor{WQ-0279-TIANDAN}{战国·田单复齐}
田单所依托的不是一支凭空出现的救国军，而是即墨守军、尚能动员的齐地城邑，以及燕军后方命令已生裂缝的时机。城防、乡里和工匠如何维持复国军的供给，史料不足以让我们画出完整的征发图，仍可由\eventanchor{WQ-0279-QIMO-LOCAL-MOBILIZATION}{战国·即墨地方动员}理解：燕军难以控制齐地基层，残城保存的组织资源便可能接到反攻上。

《史记》把火牛夜出写成复国的明亮一刻；这个场景之所以流传，正在于它把夜袭、惊扰和军心变化压缩进可见的牲畜、火光与城门。火牛装饰、夜战的每一个动作，以及由一击而收复各城的节奏，都有传奇化痕迹。故\eventanchor{WQ-0279-TIANDAN-HUONIU}{战国·即墨火牛}可以保留为后世叙事所认定的转折，不能承担整个复国过程的解释。收复城邑还需要粮食、向导、地方归附与连续的军事行动，具体年月和推进次序在文献中也有差异。

较稳妥的结论是，齐在前279年前后恢复了大部分失地，燕的占领体系随之退却。田单的胜利表明，攻城所得未必等于可以长期统治的土地；也表明残存据点一旦接上地方资源和对手的政治裂缝，仍能重新联结为国家。只是\eventref{WQ-0284-QI}{战国·燕破齐}造成的损耗不能由一次反攻抹平：齐复国之后仍有城邑与财富，却再难恢复战前那种持续向外施压的战略位置。
\end{event}

复国首先是让军需重新变成常态财政。前271年前后，襄王、田单一系须修复临淄与东部城邑的赋入，安置军民并重新接上田租、市场和道路。税额、官署和恢复速度没有连续档案，\eventanchor{WQ-0271-QI-TIANDAN-FISCAL}{战国·齐国复国财赋}与\eventanchor{WQ-0271-QI-RECOVERY-TAX}{战国·齐地税入恢复}不应被扩写为完整制度表；它们所指向的约束却明确：若没有稳定财源，复国军便无法守住西线。

同年秦客卿灶攻取刚、寿。攻取方式不详，但\eventanchor{WQ-0271-QIN-QI-GANGSHOU}{战国·秦取刚寿}说明齐的恢复并非在无人干扰的东部进行。齐方在西部重新部署守备和粮运的细节未存，\eventanchor{WQ-0271-QI-WESTERN-DEFENSE}{战国·齐西部守备}只能谨慎标作复国国家面对的必要反应。富庶仍在，战略余量却已缩小：齐必须把资源优先放在修复、守边和自保上，而很难再以灭宋时的姿态塑造中原秩序。

\subsection{最后的自保：孤立不是一个人的故事}

前265年齐襄王去世，齐王建年少继位。王室、后宫和相国网络如何分掌文书、使聘和对外信息，没有可连续追踪的齐国档案；《史记》和《战国策》关于后胜、君王后的叙述，也多由统一之后的结果回望。因而\eventanchor{WQ-0265-QI-JIAN}{战国·齐王建继位}、\eventanchor{WQ-0265-QI-WANGJIAN-REGENT}{战国·齐王建少主}与\eventanchor{WQ-0265-QI-JIAN-REGENTS}{战国·齐廷近臣}不该合写成“某一近臣完全操纵齐王”的简便解释。

更可信的结构是，复国后的继承本已脆弱，齐又面临秦在韩、赵、魏、楚、燕方向逐步推进的局面。对东部齐廷而言，谨慎不援也许能暂缓直接冲突；代价是它未能同赵、楚、燕形成可持续的共同防线。后胜受秦厚赂、劝齐不救诸侯的情节，带有秦汉叙事乐于归责近臣的形状；\eventanchor{WQ-0221-QI-HOUSHENG-ISOLATION}{战国·齐廷孤立}仍保存了不可回避的结果：齐长期未能把自己的东部资源同诸国的抗秦力量接上。

\begin{event}{公元前221年}{齐王建出临淄降秦，东方富国终于失去王号}{可靠纪年}{临淄与齐地}\eventanchor{WQ-0221-QI-JIAN-SURRENDER}{战国·齐王建降秦}
当秦先后灭韩、赵、魏、楚、燕，齐没有组织起西向救援。前221年齐王建出临淄降秦，齐成为最后覆亡的诸侯国；\eventref{WQ-0221-QI-JIAN-SURRENDER}{战国·齐王建降秦}所对应的王权终结，是可靠的年代节点。王建是否完全受后胜操纵，却是《史记·田敬仲完世家》与《战国策·齐策》等后出叙事形成的解释，不能把它当作已经录下的宫廷因果。

田氏在前386年所争得的册命、宗庙和使聘语言，至此不再属于一个独立的齐王室。临淄和东部城邑并没有从地面消失，而是被纳入秦的郡县与军政秩序；六国王权的终结也把齐长期积累的市场、工匠、道路和人力转入新的帝国框架。最后的降服因此不是“富而不能战”的道德寓言，而是外交孤立、战后恢复受限、继承与信息渠道脆弱，以及秦连续吞并共同作用的结果。
\end{event}

秦的征服也没有抹去齐地的旧名。前209年秦末失序，田儋在齐地自立为齐王，田荣等人参与组织兵众；前208年田氏支系又围绕王号和军权发生争夺。它们见于《史记·田儋列传》《陈涉世家》《项羽本纪》，亲属关系、拥立次序和早期分工均有异说，故\eventanchor{QIN-0209-TIANDAN-QI}{秦末·田儋复齐}、\eventanchor{QIN-0209-TIANRONG-QI-POLITICS}{秦末·田荣齐地集团}与\eventanchor{QIN-0208-QI-TIANJIA-STRUGGLE}{秦末·齐地复国分争}只能说明旧国名、田氏网络和城邑资源仍能成为动员语言，不能把它们当作战国齐国原封不动地复活。

\begin{sourcenote}[材料的范围]
田氏易姓与齐、燕战争，可用《史记》诸世家、列传和《资治通鉴》的编年脊柱互相校读；《孟子》的对话保存思想争论，《战国策》的说辞保存游说传统，均不宜当作逐字实录。现存传本都经过结集、编订或史家组织，读法应与《国语》一类国别叙事相同：它们特别能照亮行动者如何说明选择，却很难提供完整军令、港口账簿或占领区统计。战役城数、火牛细节、孟尝君出关和后胜受赂尤其如此。

稷下的制度形态更难凭一条材料坐实。现代研究多从《史记》《荀子》《盐铁论》等的不同成书层次、临淄城市考古与战国游士流动来讨论；《管子》相关篇章则可提出盐铁、市场与物资动员的问题，不能直接充作田齐当年法令。由此可以确认齐曾以资源吸纳知识与声望，却不能把“稷下学宫”写成制度完备、面貌不变的学院，也不能把海滨财富夸张成无需外交和组织即可支配的力量。
\end{sourcenote}

\begin{turningpoint}[制度得失：财富、外交与安全]
田氏以周命把既成的家族优势转为列国承认，齐的城邑、海滨之利和临淄网络又使它能够吸纳军队与游士。这样的富庶扩大了选择，却没有取消选择的代价：前314年入燕取得的是短期干预的机会，同时积累了燕国重建的动力；前286年灭宋把交通、市场和人口变成齐的资源，也把齐塑造成共同的威胁；前284年的联盟进军证明，财富若不能接上稳固外交和危机中的指挥体系，便会成为他国共同争夺的对象。

田单的收复揭示，占领与复国都取决于地方社会是否仍可被动员；复国后的赋入修复又提醒人们，恢复城邑不等于恢复战略位置。稷下提供的是一种文化与人才的国家能力，不能替代安全；安全一旦崩解，俸禄、辩论与市场也须重新依附于能够守住道路、城邑和盟约的政治秩序。齐亡的代价最终由王室、城邑居民、被征发的乡里与转入秦政的工匠共同承担，而不应只归结为某一位国君或近臣的一次选择。
\end{turningpoint}
